\subsection{Model Architecture}
Following prior work using Wav2vec 2.0 on stuttered speech detection, we chose another Speech SSL model,  WavLM Large~\cite{Chen2021WavLMLS} as our backbone since it significantly outperforms Wav2vec 2.0 on most SUPERB~\cite{yang21c_superb} benchmark tasks.
Our model structure is shown in Figure~\ref{fig:model_arch}.
Furthermore, we adopt the recent Hierarchical Convolution Interface~(HConv.) as our method for utilizing the WavLM backbone according to recent work~\cite{shih_interface}, which was shown to improve overtaking a learnable weighted sum of layer activations.
As pointed out in~\cite{bayerl22b_interspeech}, the information for stuttering detection exists in multiple layers, so we decided to use HConv. due to its ability to aggregate information across multiple layers.
After that, we add a few non-linear layers as the prediction head to classify individual frames as stuttered or non-stuttered.

In~\cite{bayerl22b_interspeech}, it was shown that phoneme recognition and stuttering detection tasks utilize the same set of layers in the upstream Speech SSL model.
Motivated by this, we decided to add an auxiliary Connectionist Temporal Classification~(CTC) Loss to predict the phoneme sequences.
The CTC loss is only applied in the pretraining stage where we have the ground truth phoneme sequences.

We roughly follow the training pipeline in~\cite{harvill22_interspeech}, our framework consists of two stages: Pretrain on synthetic augmentation of LibriSpeech 360 and Finetune on SEP-28K.

\noindent
\textbf{Pretrain on LibriSpeech}:
Different from previous works, we add our synthetic augmentations between the convolution layers and the transformer layers of WavLM rather than adding augmentation directly on the input waveform.
We keep the augmentation types the same as~\cite{harvill22_interspeech}, which includes artificial prolongation, and word/sound repetition.
% \david{Can you write a couple sentences describing the types of augmentations that are added?}
By doing so, we significantly save on computation and speed up the online augmentation.
Suppose the outputs of the prediction head for each input utterance are $\mathbf{o}=\left[ \mathbf{o}_1,\mathbf{o}_2,\dots \mathbf{o}_T \right]$, where $\mathbf{o}_i=\left[ o^p_i, o^n_i\right]$, corresponding to the probability of positive and negative for frame $i$ respectively.
The labels are $\mathbf{l}=\left[ \mathbf{l}_1,\mathbf{l}_2,\dots \mathbf{l}_T \right]$.
The loss is calculated as a cross entropy between the prediction head outputs and the label sequence.
\begin{align}
    \mathcal{L}_{\mathrm{dis}} = \frac{1}{T}\sum_{i=1}^{T}{
    -o^p_i \log{l^p_i}
    -o^n_i \log{l^n_i}},
\end{align}
$\mathcal{L}_{\mathrm{ctc}}$ is calculated between the frame level output of Phoneme Prediction head and the ground truth phoneme sequence in LibriSpeech~(which is obtained from a force aligner).
The overall pretrain loss is:
\begin{align}
    \mathcal{L}_{\mathrm{pretrain}} = \mathcal{L}_{\mathrm{dis}} + w_{\mathrm{ctc}} \cdot \mathcal{L}_{\mathrm{ctc}},
\end{align} where $w_{\mathrm{ctc}}$ is a scalar hyperparameter.

\noindent
\textbf{Finetune on SEP-28K}:
For the fine-tuning stage, since we do not have frame-level labels for SEP-28K, we simply take the timestep with the largest positive class prediction to calculate cross entropy with the utterance level target in SEP-28K, which is same as~\cite{harvill22_interspeech}. Formally,
\begin{align}
   t'= \mathrm{argmax}_{1\leq i \leq T}{~o^p_i},
\end{align}
\begin{align}
   \mathcal{L}_{\mathrm{dis}} = 
    -o^p_{t'} \log{l^p_{t'}}
    -o^n_{t'} \log{l^n_{t'}}.
\end{align}


% \subsection{Voice Conversion}
% To solve the data scarcity problem of stuttering detection, we came up with the idea to generate more utterance from currently available SEP-28K dataset.
% Though our preliminary experiment, we find out that voice conversion model from ElevenLabs preserve the characteristic of stuttering while changing the voice chacheristic of the speaker.
% Hence, we decide to use ElevenLabs to convert speech utterance within SEP-28K to enhance the diversity of data, hoping that it will improve the generalizability of our model.
% Practically, we generate equal amount of voice converted SEP-28K utterance for positive and negative examples in our training set.


\subsection{Implementation Details}
\textbf{Model and training details:}
% Before pretraining, we first preprocess all LibriSpeech utterances by passing them through the WavLM convolutional waveform encoder and caching them on disk.
% During the pretraining stage, we load the cached convolutional features from disk and randomly inject synthetic disfluencies onto the convolution features.
% We follow the same probability of synthetic disfluencies injection as in~\cite{harvill22_interspeech}.
The scalar hyperparameter $w_{\mathrm{ctc}}$ is set to $0.3$ empirically to balance the magnitude between the two losses.
For the finetuning stage, we feed the entire audio without augmentation and calculate a global level loss with its corresponding label.
In both stages, the HConv. interface and the stuttering detection head are trainable.
We use Adam optimizer with a learning rate of $1e-4$.
We select the checkpoint with the lowest validation loss for the pretraining stage and for the finetuneing stage, we select the one with the highest F1 score on the validation set.

\noindent
\textbf{Dataset:}
For pretraining, we use LibriSpeech 360hr and for SEP-28K, we follow the split specified in~\cite{harvill22_interspeech} for SEP-28K as they make the dataset balanced across positive and negative classes.